\SourceLeaf{75}
\Needspace{10\baselineskip}\begin{center}
\textbf{BEMERKUNGEN ZU DEN BEIDEN VORSTEHENDEN FRAGMENTEN ÜBER EINEN EULERSCHEN SATZ.}
\end{center}

\noindent\hspace*{\parindent}Der hier von \textsc{Gauss} behandelte Satz von \textsc{Euler} besagt Folgendes: »Die Funktion
$V(x, y, p, q, r, s, .\,.\,.)$ von $x, y, p=\dfrac{dy}{dx}$, $q=\dfrac{dp}{dx}$, $r=\dfrac{dq}{dx}$, $s=\dfrac{dr}{dx}$, $.\,.\,.$
liefert genau unter der Bedingung
\[
N-\dfrac{dP}{dx}+\dfrac{d^2Q}{dx^2}-\dfrac{d^3R}{dx^3}+\dfrac{d^4S}{dx^4}-.\,.\,.=0
\]
ein exaktes Differential $Vdx$, wenn $N, P, Q, R, S, .\,.\,.$ für die partiellen Ableitungen
\[
\begin{aligned}N=\dfrac{\partial V}{\partial y},\quad P=\dfrac{\partial V}{\partial p},\quad Q=\dfrac{\partial V}{\partial q},\\ R=\dfrac{\partial V}{\partial r},\quad S=\dfrac{\partial V}{\partial s},\ .\,.\,.
\end{aligned}
\]
von $V$ stehen«. \textsc{Euler} hat diesen Satz und seine Umkehrung im dritten Teil
seiner »Institutiones calculi integralis«, und zwar in \S\,92 des Anhangs über Variationsrechnung,
ausgesprochen und aus den vorangehenden Entwicklungen der Variationsrechnung bewiesen.
\textsc{Gauss} hat den oben wiedergegebenen direkten Beweis des genannten Satzes an der betreffenden
Stelle seines Handexemplars von \textsc{Euler}s Institutiones calculi integralis aufgezeichnet.
Den im zweiten Fragment entwickelten Beweis der Umkehrung hat \textsc{Gauss} in sein
Exemplar von \textsc{Klügel}s Lexicon Bd.\,I geschrieben.

\begin{flushright}
\textsc{Fricke}.
\end{flushright}

